% ===================================================================
%  TABLO No 3 — Çocukluk ve gençlik.
%  Traduction 2026 d'apres texto/io, controlee sur texto/fr et
%  texto/hi. Consigne : voir texto/tr/10-tablo-01.tex.
%
%  Deux scenes, et l'ido subdivise beaucoup plus que le francais :
%  huit des intertitres n'ont pas de vis-a-vis francais. Le turc suit
%  l'ido, comme les autres traductions.
% ===================================================================

%%K t03-apar-1 apar
\VUsaut{3.0mm}
\VUtitre{15.0pt}{60}{TABLO No 3}
\VUsaut{1.6mm}
\VUtitre{11.4pt}{0}{Çocukluk ve gençlik.}
\VUsaut{3.4mm}

%%K t03-apar-2 apar
(3 ve 4 numaralı tablolar, dört \VUgras{aile sahnesinde}
\VUgras{mevsimlerin}, \VUgras{yaşların}, \VUgras{ailenin},
\VUgras{giysilerin} ve \VUgras{sağlığın} temel sözleri öne çıksın diye
düzenlenmiş.)
\VUsaut{4.0mm}

%%K t03-tit-1 sub
\VUcentre{12.0pt}{0}{\textit{Birinci sahne.}}
\VUsaut{3.0mm}

%%K t03-tit-2 sub
\VUpk{10.2pt}{40}{Köyde bir vaftiz.}
\VUsaut{2.4mm}

%%K t03-tit-3 sub
\VUpk{10.2pt}{40}{İlkbahar.}
\VUsaut{3.0mm}

%%K t03-c1-01-1 p
1. --- \VUgras{İlkbahardayız}, \VUgras{mart}, \VUgras{nisan} ya da
\VUgras{mayıs} ayında, \VUgras{haftanın} bir gününde ---
\VUgras{pazartesi}, \VUgras{salı} ya da \VUgras{çarşamba}. Sabahın
\VUgras{saat onu}.

%%K t03-c1-02-1 p
2. --- \VUgras{Hava} güzel, \VUgras{esinti} duru. Güneş parlıyor. Birkaç küçük
\VUgras{bulut} \textsuperscript{(2)} mavi \VUgras{gökte}
\textsuperscript{(1)} yükseliyor.

%%K t03-c1-03-1 p
3. --- \VUgras{Ağaçlarda} daha \VUgras{yaprak} yok denecek kadar az. İnce
\VUgras{kavakta} \textsuperscript{(3)} ve yüksek \VUgras{çınarda}
\textsuperscript{(17)} şimdilik yalnız tomurcuk var.

%%K t03-c1-04-1 p
4. --- \VUgras{Kırlangıçlar} \textsuperscript{(4)} dönmüş, ve seve seve
cıvıldayarak şu \VUgras{külahın} \textsuperscript{(7)} çevresinde dönüyorlar,
ki o külah \VUgras{çan kulesinin} \textsuperscript{(5)}; o kule de köyün
\VUgras{kilisesinin} \textsuperscript{(6)} üstünde yükselir.

%%K t03-tit-4 sub
\VUsaut{3.4mm}
\VUpk{10.2pt}{40}{Kilise.}
\VUsaut{3.0mm}

%%K t03-c1-05-1 p
5. --- Bu kilise, büyük \VUgras{kapısı} \textsuperscript{(13)} ile büyük
\VUgras{saatiyle} \textsuperscript{(8)}, çok yalın. \VUgras{Küçük ibre}
\textsuperscript{(9)} X sayısında: \VUgras{saatleri} gösterir.
\VUgras{Büyüğü} \textsuperscript{(10)} XII (öğle) üstünde; \VUgras{dakikaları}
gösterir.

%%K t03-c1-06-1 p
6. --- Kulede \VUgras{çanlar} \textsuperscript{(11)} sevinçle çalıyor. Çatının
ucunda taştan küçük bir \VUgras{haç} \textsuperscript{(12)} yükselir.

%%K t03-c1-07-1 p
7. --- Kilisenin yalnız büyük \VUgras{kapısı} \textsuperscript{(13)} yok,
yanda küçük bir kapısı da var. \VUgras{Papaz} \textsuperscript{(15)},
\VUgras{cüppesiyle}, \VUgras{giysilikten} \textsuperscript{(14)} çıkıp
\VUgras{papaz evine} \textsuperscript{(19)} doğru o kapıdan gidiyor.

%%K t03-c1-08-1 p
8. --- Yanında işte \VUgras{sütçü kadın} \textsuperscript{(24)} ile
\VUgras{eşeği} \textsuperscript{(23)}. Çocuklara güzel sütünü ulaştırdığı
şehirden dönüyor.

%%K t03-tit-5 sub
\VUsaut{4.0mm}
\VUpk{10.2pt}{40}{Meyve bahçesi.}
\VUsaut{3.4mm}

%%K t03-c1-09-1 p
9. --- Usta Aliboron (eşek) bu sabah gezisinden çok hoşnut. Yandaki
\VUgras{bahçenin} havasını seve seve içine çekiyor; o havayı
\VUgras{çiçekler} \textsuperscript{(20)} kokutur, ki
\VUgras{meyve ağaçlarını} \textsuperscript{(18)} örterler. O
\VUgras{şeftali ağaçları} \textsuperscript{(18)} ile
\VUgras{kayısı ağaçları} \textsuperscript{(19)} bütün bütün pembe ve beyaz,
ve iyi bir ürün umdurur.

%%K t03-c1-10-1 p
10. --- Bu arada küçük \VUgras{arılar} \textsuperscript{(22)}, ki
\VUgras{kovandan} \textsuperscript{(21)} gelirler, oraya gidip vızıldayarak
tatlı \VUgras{ballarını} devşiriyor.

%%K t03-c1-11-1 p
11. --- Ama bu ne oluyor? Köyün çocukları koşarak geliyor ve ürken
\VUgras{ördekleri} \textsuperscript{(25)} dağıtıyor. Bu, kâtibin oğlunun
\VUgras{vaftizinden} sonra kiliseden çıkan alay.

%%K t03-c1-12-1 p
12. --- Küçük Jak, \VUgras{beşiğinde} \textsuperscript{(26)}, çanların sevinçli
sesiyle uyanıyor, ve alayı görmek isteyen kız kardeşi Madlen, annesine gelip
saçını eşikte taramasını ve kendisini giydirmesini söylüyor.

%%K t03-c1-13-1 p
13. --- Annesi kabul ediyor. \VUgras{Başörtüsünü} \textsuperscript{(28)},
\VUgras{korsesini} \textsuperscript{(29)} ve \VUgras{önlüğünü}
\textsuperscript{(30)} giyip kapının önünde küçük bir sandalyeye oturuyor.
Madlen'in üstünde yalnız \VUgras{iç eteği} \textsuperscript{(31)} ile
\VUgras{gömleği} \textsuperscript{(32)} var.

%%K t03-c1-14-1 p
14. --- Ne kadar sevinçli! Sevdiği çiçekleri unutuyor —
\VUgras{karanfillerini} \textsuperscript{(34)}, ki üstlerine
\VUgras{kelebekler} \textsuperscript{(35)} konar, \VUgras{lalelerini}
\textsuperscript{(36)}, \VUgras{inci çiçeklerini} \textsuperscript{(37)}, ki
\VUgras{saksılarda} \textsuperscript{(38)}, ve o saksılar
\VUgras{duvarın} \textsuperscript{(33)} üstünde, bir de öyle güzel bir çit
yapan \VUgras{güllerini} \textsuperscript{(39)}. Alaydaki hanımların değerli giysilerine bakıyor.

%%K t03-c1-15-1 p
15. --- Köyün oğlanları giysilere daha az bakar. \VUgras{Şeker}
\textsuperscript{(55)} ya da \VUgras{madenî para} \textsuperscript{(52)} almak
için öne atılıp itişiyorlar. İçlerinden biri \textsuperscript{(40)} ağlıyor.

%%K t03-c1-16-1 p
16. --- Bir başkası \VUgras{köpeğin} \textsuperscript{(42)} pençesini ezmiş;
köpek de ciyaklayarak kaçıyor ve \VUgras{tavuk} \textsuperscript{(43)} ile
\VUgras{civcivlerini} \textsuperscript{(44)} kaçırıyor.

%%K t03-tit-6 sub
\VUsaut{5.0mm}
\VUpk{10.2pt}{40}{Vaftiz alayı.}
\VUsaut{4.0mm}

%%K t03-c1-17-1 p
17. --- Alayın önünde \VUgras{sütnine} \textsuperscript{(45)} yürüyor. Başında
uzun kurdeleli güzel bir \VUgras{başlık} \textsuperscript{(46)} var.
\VUgras{Yeni doğanı} \textsuperscript{(47)} taşıyor, ki baştan aşağı
\VUgras{danteller} \textsuperscript{(48)} içinde sarılı.

%%K t03-c1-18-1 p
18. --- Sütninenin arkasından \VUgras{vaftiz babası} \textsuperscript{(49)} ile
\VUgras{vaftiz anası} \textsuperscript{(53)} geliyor. Şeker ile para
dağıtıyorlar. Vaftiz babasının elinde \VUgras{eldivenleri}
\textsuperscript{(51)}, başında \VUgras{silindir şapkası}
\textsuperscript{(50)} var. Vaftiz anası elinde güzel bir \VUgras{buket}
\textsuperscript{(54)} tutuyor.

%%K t03-c1-19-1 p
19. --- Onların arkasında çocuğun \VUgras{amcasını} \textsuperscript{(56)} ile
\VUgras{yengesini} \textsuperscript{(58)} görüyoruz. Yengenin başında şık bir
\VUgras{şapka} \textsuperscript{(59)}, amcanın üstünde siyah bir
\VUgras{redingot} \textsuperscript{(57)} ile beyaz bir yelek var. Sonra
\VUgras{kuzen} \textsuperscript{(60)} ile \VUgras{kuzini} \textsuperscript{(61)}
geliyor. Kuzin, yeni doğanın \VUgras{kız kardeşi} \textsuperscript{(62)} küçük
Berta'nın elini tutuyor.

%%K t03-c1-20-1 p
20. --- Berta'nın öteki \VUgras{erkek kardeşi} \textsuperscript{(63)} arkadan
yürüyor. Bir \VUgras{akrabaya} \textsuperscript{(64)} kolunu vermiş. Ailenin
\VUgras{dostları} \textsuperscript{(65)} onlardan sonra geliyor.

%%K t03-tit-7 sub
\VUsaut{4.6mm}
\VUpk{10.2pt}{40}{Seyirciler.}
\VUsaut{3.4mm}

%%K t03-c1-21-1 p
21. --- Alayın yanında işte bir \VUgras{dilenci} \textsuperscript{(66)},
\VUgras{koltuk değneğine} \textsuperscript{(67)} dayanmış. Sadaka istiyor.

%%K t03-c1-22-1 p
22. --- Öte yanda çok \VUgras{kibar} \textsuperscript{(68)} bir oğlan var.
Şapkasını kaldırıp tanıdıklarına \VUgras{« günaydın »} diyor.

%%K t03-c1-23-1 p
23. --- Köylüler vaftiz alayının geçişini görmek için evlerinden çıkmış.
\VUgras{Pamuklu takkesini} giymiş bir \VUgras{köylü} \textsuperscript{(69)} ile
yanında bir \VUgras{köylü kadın} \textsuperscript{(70)} görüyoruz. Sonra bir
\VUgras{yaşlı kadın} \textsuperscript{(71)}, \VUgras{değneğine}
\textsuperscript{(72)} dayanmış. Sonunda \VUgras{değirmenci}
\textsuperscript{(73)}, geniş \VUgras{keçe şapkasıyla}
\textsuperscript{(74)}, ve \VUgras{takunya} \textsuperscript{(75)} giymiş,
çocuğuna yürümeyi öğreten genç bir köylü kadın.

%%K t03-c1-24-1 p
24. --- Sahne çok canlı.
\VUsaut{4.0mm}

%%K t03-tit-8 sub
\VUcentre{12.0pt}{0}{\textit{İkinci sahne.}}
\VUsaut{3.0mm}

%%K t03-tit-9 sub
\VUpk{10.2pt}{40}{Halk şenliği.}
\VUsaut{2.4mm}

%%K t03-tit-10 sub
\VUpk{10.2pt}{40}{Su sporları ve kayık yarışı.}
\VUsaut{3.0mm}

%%K t03-c2-01-1 p
1. --- \VUgras{Yaz} gelmiş. \VUgras{Haziran} (Temmuz veya \VUgras{Ağustos})
ayının yakıcı güneşi gençleri yıkanmaya ve kayık çekmeye çağırıyor.

%%K t03-c2-02-1 p
2. --- Irmağın sağ kıyısında kayıkçılar su sporlarına ve yarışa girmek için
soyunuyor.

%%K t03-c2-03-1 p
3. --- İçlerinden biri \VUgras{potinlerini} \textsuperscript{(1)} çıkarıyor;
bağlarını (düğmelerini) çözüyor. İki arkadaşı çoktan hazır. Üstlerinde
yalnız \VUgras{fanilaları} \textsuperscript{(2)} ile \VUgras{mayoları}
\textsuperscript{(3)} var. Arkadaşlarını beklerken konuşuyorlar. Öteki
kıyıda kurulan panayır ile şenlikten söz ediyorlar.

%%K t03-c2-04-1 p
4. --- Yanlarında yerde arkadaşlarının \VUgras{giysileri} durur: bir
\VUgras{çift} \VUgras{çizme} \textsuperscript{(4)}, \VUgras{ayakkabı}
\textsuperscript{(5)}, \VUgras{çorap} \textsuperscript{(6)}, \VUgras{keçe}
\textsuperscript{(7)} ile \VUgras{hasır} \textsuperscript{(8)} şapka, bir de bir
\VUgras{boyun bağı} \textsuperscript{(9)}.

%%K t03-c2-05-1 p
5. --- Bir delikanlı \VUgras{ceketini} \textsuperscript{(10)} özenle katlamış.
Onu bir havluya koyuyor; yanındaki biraz sonra o havluya ikisinin de
\VUgras{takımlarını} bağlayacak.

%%K t03-c2-06-1 p
6. --- Altıncı oğlan geri kalmış. Başında hâlâ \VUgras{kasketi}
\textsuperscript{(11)} var. Ne \VUgras{gömleğini} \textsuperscript{(12)}
çıkarmış, ne \VUgras{yakasını} \textsuperscript{(13)}, ne de
\VUgras{kolluklarını} \textsuperscript{(14)}. Elinde \VUgras{yeleğini}
\textsuperscript{(17)} tutuyor ve hâlâ \VUgras{pantolonu}
\textsuperscript{(16)} ile \VUgras{pantolon askıları} \textsuperscript{(15)}
üstünde. Gelecek yarışa besbelli hazır olamayacak.

%%K t03-c2-07-1 p
7. --- Gençlerin yanında \VUgras{devedikenleriyle} \textsuperscript{(18)}
çevrili bir patika ırmağın kıyısına iner; orada yaşlı Jak,
\VUgras{sazların} \textsuperscript{(19)} arasında, akşam atılacak
\VUgras{havai fişekleri} \textsuperscript{(20)} hazırlıyor.

%%K t03-tit-11 sub
\VUsaut{4.4mm}
\VUpk{10.2pt}{40}{Yarış.}
\VUsaut{3.4mm}

%%K t03-c2-08-1 p
8. --- Irmakta, şemsiyelerinin gölgesinde, birkaç hanım bir \VUgras{sandalla}
\textsuperscript{(21)} çıkmış. Çok heyecanlı olan yarışı seyrediyorlar. Her
\VUgras{kayıkta} \textsuperscript{(22)} dört \VUgras{kürekçi}
\textsuperscript{(24)} var güçleriyle çekiyor; \VUgras{dümenci}
\textsuperscript{(26)} ise \VUgras{dümeni} \textsuperscript{(25)} tutup ince
kayığa yön veriyor. \VUgras{Kürekler} \textsuperscript{(23)} suya tempoyla
düşüyor, ve hafif kayıklar baş döndürücü bir hızla koşuyor.

%%K t03-c2-09-1 p
9. --- Birkaç genç, köprünün ayağı yanında, \VUgras{yağlı direk}
\textsuperscript{(28)} ödülünü kazanmaya çalışıyor. İçlerinden biri, eğilip
kayan direkte özenle ilerleyerek, ucuna bağlı küçük \VUgras{bayrağa}
\textsuperscript{(29)} varmak istiyor. Talihsiz \VUgras{rakibi}
\textsuperscript{(30)} suya düşmüş. Yüzerek kıyıya dönüyor.

%%K t03-c2-10-1 p
10. --- Büyük oğlanlar \VUgras{su cengi} \textsuperscript{(32)} oynuyor,
\VUgras{iskelenin} \textsuperscript{(33)} yanında. Büyük bir
\VUgras{kalabalık} \textsuperscript{(31)}, \VUgras{köprünün}
\textsuperscript{(27)} \VUgras{korkuluğuna} eğilmiş ve \VUgras{kıyıda} tıklım
tıklım toplanmış, bütün bu oyunları seyrediyor. Ama kimi seyirci dönüp \VUgras{tırmanma direği}
\textsuperscript{(45)} oyununa bakıyor, ya da \VUgras{ödül dağıtımına} giden
\VUgras{belediye başkanının alayını} süzüyor.

%%K t03-c2-11-1 p
11. --- En önde işte birkaç \VUgras{yumurcak} \textsuperscript{(35)},
\VUgras{bandocuların} \textsuperscript{(36)} önünden koşuyorlar; sonra \VUgras{belediye
başkanı} \textsuperscript{(37)}, yardımcıları ve küçük şehrin \VUgras{belediye
meclisi} geliyor. En arkadan \VUgras{itfaiyeciler} \textsuperscript{(38)}
yürüyor.

%%K t03-tit-12 sub
\VUsaut{5.0mm}
\VUpk{10.2pt}{40}{Panayır.}
\VUsaut{3.6mm}

%%K t03-c2-12-1 p
12. --- \VUgras{Panayır alanında} \textsuperscript{(39)} büyük bir kalabalık
var. İnsanlar türlü türlü \VUgras{çadırları} görmeye geliyor:
\VUgras{atlıkarınca} \textsuperscript{(40)}, \VUgras{sirk}
\textsuperscript{(41)}, ki orada gösteri başlamış bile; \VUgras{halk balosu}
\textsuperscript{(42)}, ki orada şehrin genç oğlanlarıyla kızları
\VUgras{dansın} tadını çıkarıyor.

%%K t03-c2-13-1 p
13. --- Dipte \VUgras{arenayı} \textsuperscript{(44)} görüyoruz; orada yiğit
ispanyol \VUgras{boğa güreşçisi} öfkeli \VUgras{boğayla}
\textsuperscript{(45)} dövüşüyor; sonra \VUgras{salıncaklı tren}
\textsuperscript{(49)} ile \VUgras{atış çadırı} \textsuperscript{(46)}, ki orada
insanlar \VUgras{tüfek} atmayı ve \VUgras{hedefi} vurmayı deniyor.

%%K t03-c2-14-1 p
14. --- Hemen yanında işte \VUgras{panayır tiyatrosu} \textsuperscript{(47)},
ki orada \VUgras{güldürüler} ile \VUgras{oyunlar} oynanır. Onun tam yanında
\VUgras{devin} \textsuperscript{(48)} ve \VUgras{cücenin} çadırı var.
Birincisinin \VUgras{boyu} iki metre on beş; ikincisi bir çocuk kadar
bile uzun değil.

%%K t03-c2-15-1 p
15. --- Sonunda işte büyük \VUgras{hayvan arabası} \textsuperscript{(50)},
\VUgras{yabanıl hayvanlarıyla} — \VUgras{pençeleri} çok güçlü olan
\VUgras{kaplanıyla} \textsuperscript{(51)}, \VUgras{yelesi} sık olan
\VUgras{aslanıyla} \textsuperscript{(52)}, iki \VUgras{zürafasıyla}
\textsuperscript{(53)} ve \VUgras{hortumunu} öyle usturuplu kullanan
\VUgras{filiyle} \textsuperscript{(54)}.

%%K t03-c2-16-1 p
16. --- Bu hayvanları eğiten \VUgras{terbiyeci} \textsuperscript{(55)}, çadırın
kapısında durur.
\VUsaut{6.0mm}
\VUfilet{22mm}
